\section{Optimization and Auxiliary Algorithm Specifications}
\label{sec:b}

This appendix formalises the mixed-integer linear programming (MILP)
layer that converts per-player point predictions into squads,
transfers, and starting lineups, together with the auxiliary
algorithms that feed or adjust those predictions: the international
and club Elo rating updates, the blended score used for live
decisions, and the rotation-risk multiplier. The three predictor
backends themselves are specified in Section~\ref{sec:methodology}
and are not repeated here. All optimisation models are built with
PuLP \citep{pulp} and solved with the bundled CBC solver \citep{cbc};
every instance in this work solves to proven optimality in well under
one second. The canonical implementation is
\texttt{optimizer/solvers.py}.

\subsection{Squad selection}
\label{sec:appendix-b-squad}

Let $\mathcal{P}$ be the candidate pool, obtained from the full
player list by removing players whose country has been eliminated
(equivalently, fixing their variables to zero). For each player
$p \in \mathcal{P}$, let $e_p$ denote the effective predicted points
summed over the planning horizon (the raw model prediction after the
scouting bonus and availability discount of
Section~\ref{sec:methodology} have been applied), $c_p$ the price in
the game's currency (millions), $\mathrm{pos}(p) \in \{\mathrm{GK},
\mathrm{DF}, \mathrm{MF}, \mathrm{FW}\}$ the position, and
$\mathrm{nat}(p)$ the country. With binary decision variables
$x_p \in \{0,1\}$ indicating squad membership, the fresh-squad
problem is
\begin{align}
\max_{x} \quad & \sum_{p \in \mathcal{P}} e_p \, x_p
  \label{eq:squad-obj}\\
\text{subject to} \quad
& \sum_{p \in \mathcal{P}} x_p = 15
  && \text{(squad size)} \label{eq:squad-size}\\
& \sum_{p :\, \mathrm{pos}(p) = r} x_p = n_r,
  \quad r \in \{\mathrm{GK}, \mathrm{DF}, \mathrm{MF}, \mathrm{FW}\}
  && \text{(position quotas)} \label{eq:squad-pos}\\
& \sum_{p \in \mathcal{P}} c_p \, x_p \le B
  && \text{(budget)} \label{eq:squad-budget}\\
& \sum_{p :\, \mathrm{nat}(p) = j} x_p \le C
  \quad \text{for every country } j
  && \text{(nationality cap)} \label{eq:squad-country}\\
& x_p \in \{0, 1\}
  \quad \forall p \in \mathcal{P}, \nonumber
\end{align}
where the position quotas are fixed by the game rules at
$(n_{\mathrm{GK}}, n_{\mathrm{DF}}, n_{\mathrm{MF}},
n_{\mathrm{FW}}) = (2, 5, 5, 3)$. The budget $B$ and cap $C$ are
stage parameters: $B = 100.0$ during the group stage and $105.0$ from
the round of 32 onward, while $C$ starts at 3 and is progressively
relaxed by the game (4 at the round of 16, 5 at the quarter-finals,
6 at the semi-finals, 8 at the final) as eliminations shrink the
pool.

\subsection{Transfers with hit accounting}
\label{sec:appendix-b-transfer}

Between rounds the game grants a stage-dependent quota of free
transfers; each transfer beyond the quota costs a fixed 3-point
deduction. Given the current squad $S_0$ and the effective free
quota $F$ (the stage allowance plus any transfers rolled over from
the previous round), the transfer problem augments
% NOTE md/code mismatch: the draft's quota constraint uses the stage
% allowance alone; solve_transfer adds rolled_over_transfers to the
% free quota (free = config.free_transfers + rolled_over_transfers).
\eqref{eq:squad-obj}--\eqref{eq:squad-country} with a continuous
overage variable $t \ge 0$:
\begin{align}
\max_{x,\, t} \quad & \sum_{p \in \mathcal{P}} e_p \, x_p - 3\, t
  \label{eq:transfer-obj}\\
\text{subject to} \quad
& \eqref{eq:squad-size}\text{--}\eqref{eq:squad-country}, \nonumber\\
& \sum_{p \notin S_0} x_p \le F + t
  && \text{(transfer quota)} \label{eq:transfer-quota}\\
& t \ge 0. \nonumber
\end{align}
Although $t$ is declared continuous, the left-hand side of
\eqref{eq:transfer-quota} and $F$ are integral and the objective is
strictly decreasing in $t$, so at any optimum
$t = \max\bigl(0, \sum_{p \notin S_0} x_p - F\bigr)$ and the
relaxation is exact. Stages with unlimited free transfers (the
initial pre-tournament selection and the reset entering the round of
32) bypass the quota machinery entirely: the fresh-squad problem of
Section~\ref{sec:appendix-b-squad} is solved and the transfer list is
reported as the set difference against $S_0$.

\subsection{Starting lineup, formation, and captaincy}
\label{sec:appendix-b-lineup}

Given a chosen 15-player squad $S$ and a target round, the lineup
solver picks the starting eleven and formation jointly. Let $s_p$ be
the player's single-round score: the effective (bonus- and
discount-adjusted) prediction when available, otherwise the raw model
prediction.
% NOTE md/code mismatch: the draft states the XI objective uses raw
% predicted_points; solve_lineup uses the effective_points column
% whenever it is present, falling back to predicted_points.
Let $\mathcal{F}$ be the set of legal formations, each fixing the
outfield counts $d_{F} = (d_F^{\mathrm{DF}}, d_F^{\mathrm{MF}},
d_F^{\mathrm{FW}})$ with one goalkeeper implied:
\[
\mathcal{F} = \{\,\text{4-4-2},\ \text{4-3-3},\ \text{4-5-1},\
\text{3-4-3},\ \text{3-5-2},\ \text{5-4-1},\ \text{5-3-2}\,\}.
\]
With binaries $y_p$ (player $p$ starts) and $f_F$ (formation $F$ is
used), the model is
\begin{align}
\max_{y,\, f} \quad & \sum_{p \in S} s_p \, y_p \\
\text{subject to} \quad
& \sum_{F \in \mathcal{F}} f_F = 1
  && \text{(one formation)} \\
& \sum_{p \in S} y_p = 11
  && \text{(eleven starters)} \\
& \sum_{p :\, \mathrm{pos}(p) = \mathrm{GK}} y_p = 1
  && \text{(one goalkeeper)} \\
& \sum_{p :\, \mathrm{pos}(p) = r} y_p
  = \sum_{F \in \mathcal{F}} d_F^{\,r} f_F,
  \quad r \in \{\mathrm{DF}, \mathrm{MF}, \mathrm{FW}\}
  && \text{(formation counts)} \\
& y_p, f_F \in \{0, 1\}. \nonumber
\end{align}
The remaining choices are made outside the MILP:
\begin{enumerate}
\item \emph{Bench order.} The four non-starters are ordered for
  automatic substitution: outfield players by descending $s_p$, with
  the backup goalkeeper always last (the game substitutes a
  goalkeeper only for the other goalkeeper).
\item \emph{Captain and vice-captain.} The solver's default is the
  starter with the highest raw predicted mean, with the runner-up as
  vice-captain. In production this placeholder is overridden by a
  ceiling-aware composite selector that interpolates between the mean
  and the 90th-percentile (q90) prediction according to the manager's
  current standings percentile, so that a manager chasing the leader
  captains for upside rather than for expectation.
\end{enumerate}

\subsection{International Elo rating update}
\label{sec:appendix-b-elo}

Country strength ratings are computed by rolling a standard Elo
recursion \citep{elo1978rating} through the full international match
history since 1872 (Section~\ref{sec:data-elo}), with
football-specific adjustments in line with common practice
\citep{hvattum2010using}. The implementation is
\url{external/international_elo.py}. Every country starts at
$R = 1500$. For each match in chronological order, with home and away
ratings $R_H$ and $R_A$, the expected home score is
\begin{equation}
\hat{s}_H = \frac{1}{1 + 10^{\,(R_A - (R_H + h))/400}},
\qquad \hat{s}_A = 1 - \hat{s}_H,
\end{equation}
where the home advantage $h$ is 60 rating points, or 0 for matches on
neutral ground. The realised score is $s_H \in \{1, \tfrac12, 0\}$
for a home win, draw, or loss ($s_A = 1 - s_H$); matches with missing
scores are treated as draws so they barely move the ratings. The
K-factor weights matches by tournament importance,
\begin{equation}
K = \begin{cases}
60 & \text{World Cup finals matches (not qualifiers)} \\
25 & \text{qualifiers} \\
10 & \text{friendlies} \\
40 & \text{all other tournaments,}
\end{cases}
\end{equation}
and a goal-difference multiplier amplifies decisive results:
\begin{equation}
g = \sqrt{\frac{m + 1}{2}}, \qquad
m = \max\bigl(1, |\text{goals}_H - \text{goals}_A|\bigr),
\end{equation}
so $g = 1$ for a one-goal margin, about $1.22$ for two, $1.41$ for
three. The update is
\begin{equation}
R_H \leftarrow R_H + K g\, (s_H - \hat{s}_H),
\qquad
R_A \leftarrow R_A + K g\, (s_A - \hat{s}_A).
\end{equation}
After the full history is processed, the per-country snapshot is the
most recent post-match rating, alongside recent-form aggregates (form
points over the last ten matches, goals for and against over the
trailing 24 months).

\subsection{Club Elo for training without lookahead}
\label{sec:appendix-b-club-elo}

Gradient-boosted model training on club data requires a club-strength
feature that was knowable before each training match. The same Elo
recursion is run over per-league match files from football-data.co.uk
with two simplifications: the K-factor is fixed at 20 (club league
schedules have no tournament-importance tiers and are less volatile
than international windows), and the home advantage of 60 is always
applied since league fixtures have a designated home side.
% NOTE md/code mismatch: the draft says the update is identical to
% the international one with missing scores scored as draws;
% compute_club_elo_history instead skips the rating update entirely
% when a score is missing (the before-match snapshot is still
% recorded).
Matches with missing scores are recorded but do not update ratings.
Crucially, the history stores each club's rating \emph{before} every
match it played. Training rows are joined to this table with a
backward as-of merge: a row for a match on date $D$ receives the
latest rating snapshot strictly before $D$. The model therefore never
sees rating information that incorporates the outcome of the match it
is being trained to predict.

\subsection{Blended score for live decisions}
\label{sec:appendix-b-blend}

Once the tournament produced realised scores, live squad decisions
were driven not by the raw model output but by a convex blend of
model prediction, realised scoring rate, and the game's own form
signal. For player $p$ at a given round,
\begin{equation}
\mathrm{final}_p = 0.55\, M_p + 0.30\, A_p
  + 0.15\, \mathrm{clip}(\phi_p, 0, 15),
\label{eq:blend}
\end{equation}
where $M_p$ is the model ensemble, the mean of the heuristic and
gradient-boosted predictions (the Poisson backend is excluded here
because it systematically over-predicted in live runs after
Matchday~2), $A_p$ is realised total points divided by matches
played, and $\phi_p$ is the form value reported by the FIFA API (a
rolling per-match average), clipped to $[0, 15]$. The weights were
hand-tuned by inspection at Matchday~3 under the principle that
realised output should dominate model priors once a player has two or
more matches of evidence; they were not learned. A principled
replacement (a Bayesian posterior with shrinkage toward the model
prior) is discussed in Section~\ref{sec:future_work}.

\subsection{Rotation-risk multiplier}
\label{sec:appendix-b-rotation}

In the final group-stage round, countries that have already clinched
qualification tend to rest starters. This is handled by a per-country
scalar $\rho_j \in (0, 1]$ applied multiplicatively to the model
ensemble before optimisation:
\begin{equation}
\tilde{M}_p = M_p \cdot \rho_{\mathrm{nat}(p)}.
\end{equation}
The multipliers were set by hand per country and matchday from squad
news and standings context, within the following bands:

\begin{center}
\begin{tabular}{lc}
\toprule
Situation & $\rho$ \\
\midrule
Clinched top spot, low-stakes fixture & 0.55--0.65 \\
Clinched, still contesting top spot   & 0.80--0.93 \\
Needs a result                        & 1.00 \\
Already eliminated                    & 0.70--0.85 \\
\bottomrule
\end{tabular}
\end{center}

A protect-list override forces $\rho = 1$ for a fixed set of eight
players contending for the tournament top-scorer award (Messi,
Mbappe, Lautaro Martinez, Ronaldo, Kane, Saka, Salah, and Haaland),
regardless of their country's multiplier, on the grounds that such
players are started even in dead rubbers. We flag this hand-set table
as a known weakness in Section~\ref{sec:lessons}; a learned
rotation-risk model is future work
(Section~\ref{sec:future_work}).
